\section*{Capítulo \ref{ch:g12:satellites-kepler} --- Satélites y movimiento planetario}

\begin{solution}{exo:g12:satellites-kepler:1}
En la plataforma: $F = GMm/R^2 = \num{3.98e17}/\num{4.06e13} \approx
\qty{9.8e3}{N}$. A $r = \qty{6.77e6}{m}$: $F \approx \qty{8.7e3}{N}$ ---
todavía el $89\,\%$ del valor en el suelo. Sus ocupantes flotan porque
\emph{caen} junto con la estación, no porque la gravedad haya
desaparecido.
\end{solution}

\begin{solution}{exo:g12:satellites-kepler:2}
$v = \sqrt{GM/R} = \sqrt{\num{3.98e14}/\num{6.37e6}} \approx
\qty{7.9e3}{m/s}$; $T = 2\pi R/v \approx \qty{5.1e3}{s} \approx
\qty{84}{min}$. Exactamente los \qty{7.9}{km/s} de la bala de cañón: una
\omterm{def:g10:universal-gravitation:satellite}{órbita} rasante \emph{es} el disparo de Newton.
\end{solution}

\begin{solution}{exo:g12:satellites-kepler:3}
$v \propto 1/\sqrt{r}$: dividida entre $\sqrt 2 \approx 1.41$.
$T \propto r^{3/2}$: multiplicado por $2^{3/2} \approx 2.8$. Ninguno de
los dos contiene $m$: no.
\end{solution}

\begin{solution}{exo:g12:satellites-kepler:4}
$v = \sqrt{\num{3.98e14}/\num{6.77e6}} \approx \qty{7.67e3}{m/s}$;
$T = 2\pi r/v \approx \qty{5.55e3}{s} \approx \qty{92}{min}$;
$\num{86400}/5546 \approx 15.6$ \omterm{def:g10:universal-gravitation:satellite}{órbitas} al día.
\end{solution}

\begin{solution}{exo:g12:satellites-kepler:5}
\omterm{def:g12:satellites-kepler:ellipse}{Elipses} con el Sol en un \omterm{def:g11:lenses-and-eye:focal}{foco}; áreas iguales en tiempos iguales;
$T^2/a^3$ idéntico para todos los planetas. Más rápido en el \omterm{def:g12:satellites-kepler:ellipse}{perihelio},
por la segunda ley (el corto segmento que barre tiene que barrer
deprisa). No: la masa del \omterm{def:g10:universal-gravitation:satellite}{satélite} se cancela --- solo entra la masa
central.
\end{solution}

\begin{solution}{exo:g12:satellites-kepler:6}
$T = \qty{2.36e6}{s}$: $M = 4\pi^2 r^3/(GT^2) = \num{2.24e27}/371
\approx \qty{6.0e24}{kg}$ --- a menos de un $1\,\%$ de
\qty{5.97e24}{kg}.
\end{solution}

\begin{solution}{exo:g12:satellites-kepler:7}
La gravedad apunta al centro de la Tierra y, en un movimiento circular,
la \omterm{def:g10:inertia:force}{fuerza} resultante apunta al centro de la circunferencia: los dos
centros coinciden. Una circunferencia sobre el paralelo de París está
centrada en el eje, no en el centro terrestre --- imposible. Quedarse
colgado solo funciona allí donde el paralelo es un círculo máximo
centrado en el centro: sobre el \emph{ecuador}.
\end{solution}

\begin{solution}{exo:g12:satellites-kepler:8}
$T = \qty{2.754e4}{s}$: $M = 4\pi^2 r^3/(GT^2) =
\num{3.26e22}/\num{5.06e-2} \approx \qty{6.4e23}{kg}$ --- el Marte de la
tabla de datos.
\end{solution}

\begin{solution}{exo:g12:satellites-kepler:9}
Tierra: $T^2/a^3 = \num{9.96e14}/\num{3.35e33} =
\qty{2.98e-19}{s^2/m^3}$. Marte: $T = \qty{5.94e7}{s}$,
$\num{3.52e15}/\num{1.18e34} = \qty{2.98e-19}{s^2/m^3}$ --- iguales,
como exige la tercera ley. Entonces
$M = 4\pi^2/(G \times \num{2.98e-19}) \approx \qty{1.99e30}{kg}$.
\end{solution}

\begin{solution}{exo:g12:satellites-kepler:10}
$v = 2\pi r/T = \qty{1.02e3}{m/s}$;
$a = v^2/r = \qty{2.72e-3}{m/s^2}$;
$g/60^2 = 9.81/3600 = \qty{2.73e-3}{m/s^2}$. La Luna está a $60$ radios
terrestres y su caída es $60^2$ veces más débil que la de la manzana:
exactamente el inverso del cuadrado --- una sola ley para las dos.
\end{solution}

\begin{solution}{exo:g12:satellites-kepler:11}
Ceres: $T = 2.77^{3/2} = \sqrt{21.3} \approx \qty{4.6}{años}$. Halley:
$T = 17.8^{3/2} \approx \qty{75}{años}$, así que el próximo \omterm{def:g12:satellites-kepler:ellipse}{perihelio} cae
hacia $1986 + 75 \approx 2061$.
\end{solution}

\begin{solution}{exo:g12:satellites-kepler:12}
A \qty{400}{km}: $v \approx \qty{7.67e3}{m/s}$; a \qty{350}{km}
($r = \qty{6.72e6}{m}$): $v \approx \qty{7.70e3}{m/s}$. El \omterm{def:g10:inertia:inventory}{rozamiento} sí
drena \omterm{def:g11:mechanical-energy:em}{energía mecánica} --- la \omterm{def:g10:universal-gravitation:satellite}{órbita} se encoge; pero durante la bajada
el \omterm{def:g11:work-of-force:work}{trabajo} de la gravedad devuelve con creces la rapidez que se lleva el
aire: $E_m$ baja mientras $E_k$ sube.
\end{solution}

\begin{solution}{exo:g12:satellites-kepler:13}
$T = \qty{3.024e5}{s}$: $M = 4\pi^2 r^3/(GT^2) = \num{1.65e31}/6.10
\approx \qty{2.7e30}{kg} \approx 1.4$ soles. La masa \emph{del planeta}
queda fuera de \omterm{def:g12:projectile-motion:range}{alcance}: se canceló en las ecuaciones de la \omterm{def:g10:universal-gravitation:satellite}{órbita}.
\end{solution}

\begin{solution}{exo:g12:satellites-kepler:14}
$T = \qty{8.862e4}{s}$, $GM = \qty{4.28e13}{m^3/s^2}$:
$r^3 = GMT^2/(4\pi^2) = \num{8.52e21}$, luego
$r \approx \qty{2.04e7}{m}$ y $h = r - R \approx \qty{1.7e7}{m} \approx
\qty{17000}{km}$ sobre el suelo marciano.
\end{solution}

\begin{solution}{exo:g12:satellites-kepler:15}
$v_a = r_p v_p / r_a = \num{8.77e10} \times \num{5.45e4} /
\num{5.25e12} \approx \qty{9.1e2}{m/s}$. Como $r_a/r_p \approx 60$, el
cometa está $60$ veces más lejos y va $60$ veces más lento: se pasa
décadas remoloneando en la oscuridad y semanas esprintando junto al Sol.
\end{solution}

\begin{solution}{pb:g12:satellites-kepler:1}
\textbf{1.} $T = \qty{86164}{s}$: el suelo gira con la Tierra
\emph{respecto de las estrellas} una vez por \omterm{def:g12:satellites-kepler:geo}{día sidéreo}; el día de
\qty{24}{h} añade el grado diario que la Tierra avanza alrededor del
Sol.

\textbf{2.} La \omterm{def:g10:inertia:force}{fuerza} resultante de un movimiento circular apunta al
centro de la circunferencia; la única \omterm{def:g10:inertia:force}{fuerza}, la gravedad, apunta al
centro de la Tierra: los dos centros son el mismo punto.

\textbf{3.} Una circunferencia suspendida sobre París estaría centrada
en el eje, al norte del centro --- excluido por el apartado 2. La \omterm{def:g10:universal-gravitation:satellite}{órbita}
debe estar en el \emph{plano ecuatorial}.

\textbf{4.} Componente normal de la \omterm{thm:g12:newtons-laws:second}{segunda ley de Newton}:
$m v^2/r = GmM/r^2$, luego $v = \sqrt{GM/r}$ ($m$ se cancela).

\textbf{5.} $T = 2\pi r/v = 2\pi\sqrt{r^3/(GM)}$; elevando al cuadrado,
$T^2/r^3 = 4\pi^2/(GM)$.

\textbf{6.} $r = \left(GMT^2/4\pi^2\right)^{1/3} = (\num{7.49e22})^{1/3}
\approx \qty{4.22e7}{m}$.

\textbf{7.} $h = r - R = \num{4.22e7} - \num{6.37e6} \approx
\qty{3.58e7}{m} \approx \qty{35800}{km}$.

\textbf{8.} $v = 2\pi \times \num{4.22e7}/\num{86164} \approx
\qty{3.07e3}{m/s}$.

\textbf{9.} $\sqrt{\num{3.98e14}/\num{4.22e7}} \approx
\qty{3.07e3}{m/s}$ --- coherente.

\textbf{10.} Radios idénticos: la masa se canceló en el apartado 4. El
\omterm{def:g10:universal-gravitation:satellite}{satélite} más pesado solo cambia la factura del lanzamiento --- más
combustible para la misma \omterm{def:g10:universal-gravitation:satellite}{órbita}.

\textbf{11.} $v = \sqrt{\num{3.98e14}/\num{6.77e6}} \approx
\qty{7.67e3}{m/s}$.

\textbf{12.} $T = 2\pi r/v \approx \qty{5.55e3}{s} \approx
\qty{92.4}{min}$.

\textbf{13.} $\num{86400}/5546 \approx 15.6$ \omterm{def:g10:universal-gravitation:satellite}{órbitas}: unos $16$
amaneceres al día.

\textbf{14.} $r_{\text{geo}}/r_{\text{EEI}} = \num{4.22e7}/\num{6.77e6}
= 6.23$; $6.23^{3/2} \approx 15.5 = \num{86164}/5546$ --- la tercera ley
de Kepler, verificada sobre la marcha.

\textbf{15.} La estación va más deprisa (\qty{7.67}{km/s} frente a
\qty{3.07}{km/s}) y, aun así, el \omterm{def:g12:satellites-kepler:geo}{geoestacionario} costó más \omterm{def:g10:energy-conservation:energy}{energía} por
\omterm{def:g10:orders-of-magnitude:unit}{kilogramo}: la subida en \omterm{def:g10:energy-conservation:energy}{energía} potencial pesa más que la pérdida de
rapidez.

\textbf{16.} $M_J = 4\pi^2 r^3/(G T^2)$.

\textbf{17.} $T = \qty{1.529e5}{s}$: $M_J = \num{2.97e27}/1.56 \approx
\qty{1.90e27}{kg}$.

\textbf{18.} $\num{1.90e27}/\num{5.97e24} \approx 318$ Tierras --- y
alrededor de $1/1000$ del Sol.

\textbf{19.} $T_J = \qty{3.576e4}{s}$: $r^3 = GM_J T_J^2/(4\pi^2) =
\num{4.11e24}$, $r \approx \qty{1.60e8}{m}$ --- Ío, a \qty{4.22e8}{m},
vuela $2.6$ veces más alto y va derivando por el cielo de Júpiter.

\textbf{20.} Aparca el repetidor en la circunferencia ecuatorial de
\qty{35800}{km} de altura, donde viaja a \qty{3.07}{km/s} y no se sale
nunca de la puntería de la antena; y esa misma ley, alimentada con el
mes de Ío, ha pesado Júpiter en $\qty{1.90e27}{kg}$ --- $318$ Tierras
--- gratis.
\end{solution}
